\paragraph{A clock read from the changing field configuration.}
A repair count and a time coordinate are different objects. The same
autonomous action does, however, supply a dynamical duration functional.
This is an application of the classical Jacobi--Maupertuis construction
\cite{Gryb2010}, not a new quantization rule. Write a natural mechanical
Lagrangian as \(L=\tfrac12G_q(\dot q,\dot q)-V(q)\), using the complete
coupled kinetic metric, or its Gauss-reduced metric on an admissible
horizontal path. Fix an energy \(E\) and an oriented regular path in the
open region \(V<E\). Define
\begin{equation}
 d\tau=\sqrt{\frac{G_q(dq,dq)}{2(E-V(q))}},\qquad
 J_E[q]=\int\sqrt{2(E-V(q))G_q(q',q')}\,ds.
 \label{eq:whitney-ephemeris-clock}
\end{equation}

\begin{proposition}[Dynamical duration on a nonturning path]
\label{prop:whitney-ephemeris-clock}
For smooth positive-definite \(G\) and smooth \(V\), the duration above
is invariant under orientation-preserving regular reparameterization.
A regular stationary path of \(J_E\), parameterized by \(\tau\), solves
the Euler--Lagrange equations of \(L\) with energy \(E\). Conversely,
every solution of that energy with positive kinetic energy is such a path,
and its \(\tau\) increments equal its original action-time increments.
For fixed \(G,V,E\), positive energy-compatible timing of that oriented
path is unique up to its origin.
\end{proposition}
\begin{proof}
Set \(T_s=G_q(q',q')/2\) and \(N=\sqrt{T_s/(E-V)}\).
Positive reparameterization multiplies both \(q'\) and \(N\) by
the same positive derivative, proving duration invariance. The lapse action
\[
 I[q,N]=\int\bigl(T_s/N+N(E-V)\bigr)\,ds
\]
has lapse equation \(T_s/N^2=E-V\). Its unique positive solution gives
\(I[q,N(q,q')]=J_E[q]\). Substitution preserves the configuration
Euler equation because the lapse variation vanishes. With
\(d\tau=N\,ds\), that equation is
\[
 \frac{d}{d\tau}(G_{ij}\dot q^j)
 -\frac12\partial_iG_{jk}\dot q^j\dot q^k+\partial_iV=0.
\]
The lapse equation is the required energy identity. Conversely this
identity and the Euler equation give the two equations for \(I\).
Uniqueness follows directly from its positive square root. The assertion
does not extend through a turning point where numerator and denominator
both vanish, or to a constant equilibrium path.
\end{proof}

This construction is covariant under changes of configuration coordinates.
For gauge independence one must use the reduced metric or keep the scalar
potential in the full covariant tangent. Arbitrary time-dependent gauge
changes followed by discarding the scalar potential do not preserve the
kinetic energy. The numerical construction below uses the specified
temporal-gauge representative throughout. Nor does a Jacobi metric on
configuration space itself define a Lorentzian spacetime metric.

There is also a useful approximation statement. Let a fixed \(C^3\)
regular path remain in a compact nonturning region, with \(G,V\) smooth
there. Its polygonal interpolants at parameter spacing at most \(\delta\)
have duration error \(O(\delta^2)\). To see this, write the duration
integrand as \(F(q,v)=\sqrt{G_q(v,v)/(2(E-V(q)))}\). On each short
segment, the configuration error is \(O(h^2)\), the velocity error is
\(O(h)\), and the latter has zero integral. Expand \(F\) to first
order: freezing its velocity derivative at the midpoint cancels that
integral, and the derivative's variation and quadratic remainder both
integrate to \(O(h^3)\). Summing gives the assertion. Smoothness is used
away from zero velocity and \(E-V=0\). This bound compares exact curve
integrals; floating-point quadrature and errors in the supplied samples
require separate control.

The executable consumer reads only the ordered decoded configuration
coordinates from a bounded self-reading software episode. Local patch
states, port responses, retained records and feedback restoration belong
to that episode's independent provenance check. The clock consumer uses
neither the recorded velocities, action timestamps nor repair counts.
It retains the same complete scalar-basis derivative in its kinetic form
and supplies the initial energy
\(E=1969/160+(493/120)\sqrt5\). Piecewise linear interpolation gives
a duration on each configuration segment. A separate implementation
integrates unrestricted tetrahedral fields instead of the producer's
one-dimensional symmetric formulas; a nonlinear reparameterization of
each segment checks duration invariance. The producer, verifier and
data are \path{code/electromagnetism/whitney_ephemeris_clock.py},
\path{code/electromagnetism/verify_whitney_ephemeris_clock.py} and
\path{code/electromagnetism/runtime/whitney_ephemeris_clock_receipt.json}.

The clock is internal to the supplied autonomous model: its rate is
reconstructed from change and the action-energy convention. Selecting that
action, energy and physical preparation from observer histories, assigning
physical units, and comparing this duration with laboratory clocks remain
additional tasks. It neither constructs a quantum clock observable nor
identifies an arbitrary event ordering with proper time.
